% Appendix: Core Code

\subsection*{Core Physics Functions}

\begin{lstlisting}[language=Python, caption={Transfer Matrix Method with Airy Summation}]
def transfer_matrix_single_layer(n_layer, d_um, wavelength_um, theta_rad,
                                  n_in=1.0, n_out=1.0):
    """
    Single-layer reflectance/transmittance via Airy summation.
    Structure: n_in | n_layer(d) | n_out
    """
    cos_in = cos(theta_rad)
    cos_layer = snell(n_in, n_layer, theta_rad)
    cos_out = snell(n_in, n_out, theta_rad)

    # Phase factor: beta = 2*pi*n_tilde*d*cos(theta) / lambda
    beta = 2.0 * pi * n_layer * d_um / wavelength_um * cos_layer

    R_total, T_total = 0.0, 0.0
    for pol in ['s', 'p']:
        # Fresnel coefficients at interfaces 0->1 and 1->2
        if pol == 's':
            r01 = (n_in*cos_in - n_layer*cos_layer) / \
                  (n_in*cos_in + n_layer*cos_layer)
            r12 = (n_layer*cos_layer - n_out*cos_out) / \
                  (n_layer*cos_layer + n_out*cos_out)
            t01 = 2.0*n_in*cos_in / (n_in*cos_in + n_layer*cos_layer)
            t12 = 2.0*n_layer*cos_layer / \
                  (n_layer*cos_layer + n_out*cos_out)
        else:
            r01 = (n_layer*cos_in - n_in*cos_layer) / \
                  (n_layer*cos_in + n_in*cos_layer)
            r12 = (n_out*cos_layer - n_layer*cos_out) / \
                  (n_out*cos_layer + n_layer*cos_out)
            t01 = 2.0*n_in*cos_in / (n_layer*cos_in + n_in*cos_layer)
            t12 = 2.0*n_layer*cos_layer / \
                  (n_out*cos_layer + n_layer*cos_out)

        exp_2ib = np.exp(2j * beta)  # Decays for absorbing media
        denom = 1.0 + r01 * r12 * exp_2ib

        r_total = (r01 + r12 * exp_2ib) / denom
        t_total = (t01 * t12 * np.exp(1j * beta)) / denom

        R = abs(r_total)**2
        # Transmission with impedance correction
        if pol == 's':
            T_pol = (n_out.real*cos_out.real) / \
                    (n_in.real*cos_in.real) * abs(t_total)**2
        else:
            T_pol = (n_out.real/cos_out.real) / \
                    (n_in.real/cos_in.real) * abs(t_total)**2

        R_total += R; T_total += T_pol

    return np.clip(R_total/2, 0, 1), np.clip(T_total/2, 0, 1)
\end{lstlisting}

\begin{lstlisting}[language=Python, caption={Planck Blackbody Spectral Radiance}]
def planck_spectral_radiance(wavelength_um, T):
    """
    I_BB(lambda,T) = 2hc^2/lambda^5 * 1/(exp(hc/(lambda*k_B*T)) - 1)
    Returns: W/(m^2.um.sr)
    """
    wl_m = wavelength_um * 1e-6
    exponent = h * c / (wl_m * kB * T)
    with np.errstate(over='ignore'):
        radiance_per_m = 2.0*h*c**2 / (wl_m**5) / (np.exp(exponent)-1.0)
        radiance = radiance_per_m * 1e-6  # Convert to per-micrometer
    radiance = np.where(np.isfinite(radiance), radiance, 0.0)
    return radiance
\end{lstlisting}

\begin{lstlisting}[language=Python, caption={Net Cooling Power Calculation}]
def compute_net_cooling_power(d_um, wavelengths_um, n_arr, k_arr,
                              T, T_amb, h_c=6.9, n_sub=1.0+0j,
                              theta_sun=np.radians(48.2), n_angles=12):
    """P_cool = P_rad - P_atm - P_sun - P_nonrad"""
    # Setup Gauss-Legendre angular quadrature
    x, w = leggauss(n_angles)
    theta_arr = (x + 1.0) * pi / 4.0
    weights = w * pi / 4.0

    # P_rad: Hemispherical integration of Planck emission
    P_rad = 0.0
    for theta, wt in zip(theta_arr, weights):
        emis = compute_emissivity_spectrum(d_um, wavelengths_um,
                                           theta, n_arr, k_arr, n_sub)
        I_bb = planck_spectral_radiance(wavelengths_um, T)
        P_rad += 2*pi * sin(theta)*cos(theta) * wt * \
                 simpson(I_bb * emis, wavelengths_um)

    # P_atm: Atmospheric counter-radiation
    t_atm = build_atmospheric_transmittance(wavelengths_um)
    P_atm = 0.0
    for theta, wt in zip(theta_arr, weights):
        emis = compute_emissivity_spectrum(d_um, wavelengths_um,
                                           theta, n_arr, k_arr, n_sub)
        eps_atm = 1.0 - t_atm ** (1.0/max(cos(theta), 1e-6))
        I_bb_atm = planck_spectral_radiance(wavelengths_um, T_amb)
        P_atm += 2*pi * sin(theta)*cos(theta) * wt * \
                 simpson(I_bb_atm * emis * eps_atm, wavelengths_um)

    # P_sun: Solar absorption (AM1.5)
    emis_sun = compute_emissivity_spectrum(d_um, wavelengths_um,
                                           theta_sun, n_arr, k_arr, n_sub)
    I_sun = build_am15_spectrum(wavelengths_um)
    P_sun = simpson(I_sun * emis_sun, wavelengths_um)

    # P_nonrad: Convection/conduction
    P_nonrad = h_c * (T_amb - T)

    P_cool = P_rad - P_atm - P_sun - P_nonrad
    return {'P_rad': P_rad, 'P_atm': P_atm, 'P_sun': P_sun,
            'P_nonrad': P_nonrad, 'P_cool': P_cool}
\end{lstlisting}

Full code (problem1.py -- problem4.py, utils.py) is available in the \texttt{code/} directory. The implementation uses NumPy for vectorized computation, SciPy for numerical integration (Gauss-Legendre quadrature and Simpson's rule), and Matplotlib for visualization. All figures are generated as PDF vector graphics suitable for publication.
